% GENERATED FILE. Edit canonical lessons, then run npm run generate:books.
% canonical-source-hash: fnv1a64:3a9f66e0280f5612
% canonical-lessons: MR-C191-vyavasthapak, MR-C191-saciv, MR-C191-hiseb, MR-C191-tapsil, MR-C191-mahiti

\chapter{Manager, Secretary, Accounts, Details, Information}
\label{ch:mr-manager-secretary-accounts-details-information}
\hlchaptermodality{191}

\begin{chapteropening}
\textbf{By the end of this chapter:} \emph{I can say a manager, a secretary, accounts, details and information.}

Complete the last lesson of chapter 191: I can say a manager, a secretary, accounts, details and information.
\end{chapteropening}

\section[\texorpdfstring{\mr{व्यवस्थापक} (vyavasthāpak)}{व्यवस्थापक (vyavasthāpak)}]{\mr{व्यवस्थापक} (vyavasthāpak) — a manager}

\label{lesson:MR-C191-vyavasthapak}



\begin{quote}
Before the new one: say the Marathi for a queue, then the Marathi for a crowd.
\end{quote}

\subsection*{You'll want to know: \mr{व्यवस्थापक}}
\textbf{\mr{व्यवस्थापक}} — \emph{vyavasthāpak} — \textquotedblleft{}a manager\textquotedblright{}.

\begin{cousinweb}
\textbf{\mr{व्यवस्थापक}} (\emph{vyavasthāpak}) is Sanskrit \emph{vyavasthāpaka}, \textquotedblleft{}one who arranges, settles\textquotedblright{}, a tatsama built on \emph{vyavasthā}, \textquotedblleft{}arrangement\textquotedblright{}. CDIAL has no entry for it.
\end{cousinweb}

\begin{grammarlens}[title={What you've built}]
One more word for this chapter.
\end{grammarlens}

\subsection*{Your turn}
\begin{itemize}
\raggedright
  \item \emph{Say it:} \emph{vyavasthāpak}
  \item \emph{Say it:} \emph{vyavasthāpak}, once more
  \item \emph{Say it:} say \emph{vyavasthāpak}
  \item \emph{Recall:} say the Marathi for a queue, then the Marathi for a crowd, then say \emph{vyavasthāpak} again
\end{itemize}

\begin{tcolorbox}[breakable,colback=teal!4,colframe=teal!35!black,title={Before you move on}]
What does \mr{व्यवस्थापक} mean? (\textquotedblleft{}A manager\textquotedblright{}.) Say it once more.
\end{tcolorbox}
\section[\texorpdfstring{\mr{सचिव} (saciv)}{सचिव (saciv)}]{\mr{सचिव} (saciv) — a secretary}

\label{lesson:MR-C191-saciv}



\begin{quote}
Before the new one: say the Marathi for a crowd, then the Marathi for a manager.
\end{quote}

\subsection*{You'll want to know: \mr{सचिव}}
\textbf{\mr{सचिव}} — \emph{saciv} — \textquotedblleft{}a secretary\textquotedblright{}.

\begin{cousinweb}
\textbf{\mr{सचिव}} (\emph{saciv}) is Sanskrit \emph{saciva}, \textquotedblleft{}companion, counsellor, minister\textquotedblright{}, taken over whole, a tatsama. CDIAL has no entry for it.
\end{cousinweb}

\begin{grammarlens}[title={What you've built}]
One more word for this chapter.
\end{grammarlens}

\subsection*{Your turn}
\begin{itemize}
\raggedright
  \item \emph{Say it:} \emph{saciv}
  \item \emph{Say it:} \emph{saciv}, once more
  \item \emph{Say it:} say \emph{saciv}
  \item \emph{Recall:} say the Marathi for a crowd, then the Marathi for a manager, then say \emph{saciv} again
\end{itemize}

\begin{tcolorbox}[breakable,colback=teal!4,colframe=teal!35!black,title={Before you move on}]
What does \mr{सचिव} mean? (\textquotedblleft{}A secretary\textquotedblright{}.) Say it once more.
\end{tcolorbox}
\section[\texorpdfstring{\mr{हिशेब} (hiśeb)}{हिशेब (hiśeb)}]{\mr{हिशेब} (hiśeb) — accounts, a reckoning}

\label{lesson:MR-C191-hiseb}



\begin{quote}
Before the new one: say the Marathi for a manager, then the Marathi for a secretary.
\end{quote}

\subsection*{You'll want to know: \mr{हिशेब}}
\textbf{\mr{हिशेब}} — \emph{hiśeb} — \textquotedblleft{}accounts, a reckoning\textquotedblright{}.

\begin{cousinweb}
\textbf{\mr{हिशेब}} (\emph{hiśeb}) is a loan from Arabic \emph{ḥisāb}, \textquotedblleft{}reckoning, account\textquotedblright{}, which reached Marathi through Persian. Its Arabic root is \emph{ḥ-s-b} (Lane).
\end{cousinweb}

\begin{grammarlens}[title={What you've built}]
One more word for this chapter.
\end{grammarlens}

\subsection*{Your turn}
\begin{itemize}
\raggedright
  \item \emph{Say it:} \emph{hiśeb}
  \item \emph{Say it:} \emph{hiśeb}, once more
  \item \emph{Say it:} say \emph{hiśeb}
  \item \emph{Recall:} say the Marathi for a manager, then the Marathi for a secretary, then say \emph{hiśeb} again
\end{itemize}

\begin{tcolorbox}[breakable,colback=teal!4,colframe=teal!35!black,title={Before you move on}]
What does \mr{हिशेब} mean? (\textquotedblleft{}Accounts, a reckoning\textquotedblright{}.) Say it once more.
\end{tcolorbox}
\section[\texorpdfstring{\mr{तपशील} (tapśīl)}{तपशील (tapśīl)}]{\mr{तपशील} (tapśīl) — details}

\label{lesson:MR-C191-tapsil}



\begin{quote}
Before the new one: say the Marathi for a secretary, then the Marathi for accounts.
\end{quote}

\subsection*{You'll want to know: \mr{तपशील}}
\textbf{\mr{तपशील}} — \emph{tapśīl} — \textquotedblleft{}details\textquotedblright{}.

\begin{cousinweb}
\textbf{\mr{तपशील}} (\emph{tapśīl}) is a loan from Arabic \emph{tafṣīl}, \textquotedblleft{}a setting out in detail\textquotedblright{}, which reached Marathi through Persian. Its Arabic root is \emph{f-ṣ-l} (Lane).
\end{cousinweb}

\begin{grammarlens}[title={What you've built}]
One more word for this chapter.
\end{grammarlens}

\subsection*{Your turn}
\begin{itemize}
\raggedright
  \item \emph{Say it:} \emph{tapśīl}
  \item \emph{Say it:} \emph{tapśīl}, once more
  \item \emph{Say it:} say \emph{tapśīl}
  \item \emph{Recall:} say the Marathi for a secretary, then the Marathi for accounts, then say \emph{tapśīl} again
\end{itemize}

\begin{tcolorbox}[breakable,colback=teal!4,colframe=teal!35!black,title={Before you move on}]
What does \mr{तपशील} mean? (\textquotedblleft{}Details\textquotedblright{}.) Say it once more.
\end{tcolorbox}
\section[\texorpdfstring{\mr{माहिती} (māhitī)}{माहिती (māhitī)}]{\mr{माहिती} (māhitī) — information}

\label{lesson:MR-C191-mahiti}



\begin{quote}
Before the new one: say the Marathi for accounts, then the Marathi for details.
\end{quote}

\subsection*{You'll want to know: \mr{माहिती}}
\textbf{\mr{माहिती}} — \emph{māhitī} — \textquotedblleft{}information\textquotedblright{}.

\begin{cousinweb}
\textbf{\mr{माहिती}} (\emph{māhitī}) is built on Marathi \emph{māhīt}, \textquotedblleft{}known\textquotedblright{}. \emph{Māhīt} is often traced, through Persian, to Arabic \emph{māhiyya}, \textquotedblleft{}what a thing is, its nature\textquotedblright{}, but no source here settles it.
\end{cousinweb}

\begin{grammarlens}[title={What you've built}]
One more word for this chapter.
\end{grammarlens}

\subsection*{Your turn}
\begin{itemize}
\raggedright
  \item \emph{Say it:} \emph{māhitī}
  \item \emph{Say it:} \emph{māhitī}, once more
  \item \emph{Say it:} say \emph{māhitī}
  \item \emph{Recall:} say the Marathi for accounts, then the Marathi for details, then say \emph{māhitī} again
\end{itemize}

\begin{tcolorbox}[breakable,colback=teal!4,colframe=teal!35!black,title={Before you move on}]
What does \mr{माहिती} mean? (\textquotedblleft{}Information\textquotedblright{}.) Say it once more.
\end{tcolorbox}
